\documentclass{article}
\usepackage{amsmath}
\begin{document}
\section*{Joint service over a finite window of actually eligible cuts}
A prepared old diagnostic remains qualified across two provisional
registrations. The question is whether it can still be delivered by
its original cut, while meeting legally required registrations and
without losing either observer's later service. At a serialized
effect cut a member whose pose source is not ready prevents the
pinned native broker from popping the queue, irrespective of the
diagnostic's own source and destination. No external atomic
preflight API is attributed to the stock broker.

Let \(W\) be the finite set of opportunities at which the old send
would otherwise be lawful and fit the actual resource profiles.
This factors out membership readiness only, not authority, link
qualification, resource use, due time or prior work. These are the
same prepared send and eligibility rules as in the pinned
effect-cut checker. Exports of two observer availability
contexts must be interpreted on these same times. If a selected
set of early registrations is \(J\), and observer \(i\) has
unready set \(M_i \subseteq W\), the still-available cuts are
\[
             W\setminus\bigcup_{i\in J}M_i.
\]
In the replay the controller checks readiness and sends at its
first available cut. It need not know which pattern of outages will
occur. Each failed preflight leaves this send's resource credits
unchanged; the specified ambient work does not consume its TX or
RX footprint on \(W\).

\paragraph{Sharp summary using caps alone.}
Suppose each \(M_i\) has size at most \(b_i\). The minimum remaining
count over all combinations described only by these per-observer
caps is \(\max(0,|W|-\sum_i b_i)\). The lower bound follows by
removing at most \(b_i\) currently surviving cuts each time.
For sharpness, order \(W\) and allocate the next unused cuts,
up to \(b_i\) to each observer until either cuts or caps run out.
Each such schedule is in the declared product support. Lean
checks the finite-set lower bound and the exact arithmetic of
that allocation; the construction is not a claim that these
outages are physically independent. A separately verified
alignment restriction can strengthen the guarantee. Raw-tick
outage counts are not the declared certificates: they must be valid
as bounds on disruptions of \(W\). Even when a raw count gives a
sound loose bound, it does not identify which qualifying cuts
are lost or certify the same maximum number of entrants.

Thus a maximal number of early registrations can be chosen by
enumerating the authorized subsets, retaining the required ones,
and enforcing \(\sum_{i\in J} b_i<|W|\). When no subset is
feasible, the old guarantee and all required pre-deadline
registration obligations are incompatible at this information and
control interface. The replay retains those required registrations
in this branch rather than treating an impossibility as authority
to delay them. Deferral is allowed only where provisionally staged
agents have no obligation that requires being registered earlier.
An equally informed direct enumerator agrees for all 1,500 checked
mask, cap and required-subset combinations.

\paragraph{Native discriminator.}
At cuts \(0.5,1.25,1.375,1.5\) seconds, the no-debt fixture has
four eligible cuts. An already accepted 4-byte receive use at
\(0.25\) seconds blocks the 3-byte old send in the first two
cuts under a 6-byte cap. Age exactly one second is still charged.
Now \(W\) has only two cuts. Each of two observers can miss a
single such cut and still leave one if acting alone. If both
must be registered early, equal marginal miss counts give
opposite outcomes: aligned misses give service, opposed misses
cover \(W\). Deferring one optional registration preserves an
old delivery in every capped history where that deferral is
allowed and needed. Both observers still execute their own native
later send at \(1.625\) seconds in all 676 runs. Of the 144
forced early cases without a count-only guarantee, 32 actually
miss the old duty; aligned controls succeed. Native forwarding
runs with the actual prior debt and route/pose callbacks.
An unguarded comparison directly checks that failed readiness
updates earlier native pose stamps, yet leaves the queue and
that send's accounting unchanged. It is not a stutter on the
full native state.

This is a constructive finite service lift of the actual one-cut
interface in a specified product family, plus a witnessed limit
of separately successful registrations. It does not assume
preventable failure, arbitrary authority to delay a participant,
state-independent eligibility in other missions, service under
unbounded outages, stock interface changes, or independently
selected mission acceptance.
\end{document}
